\section*{Capítulo \ref{ch:b2:ellingham} --- Diagramas de Ellingham e metalurgia}

\begin{solution}{exo:b2:ellingham:1}
\ce{4Cu + O2 -> 2Cu2O}; \ce{4/3Al + O2 -> 2/3Al2O3}; \ce{2C + O2 -> 2CO};
\ce{C + O2 -> CO2}.
\end{solution}

\begin{solution}{exo:b2:ellingham:2}
$\Delta_r H^\circ = 2(-350.46) = \qty{-700.92}{kJ}$; $\Delta_r S^\circ = 2(43.65)
- 2(41.63) - 205.15 = \qty{-201.1}{J/K}$. Logo $\Delta_r G^\circ = -700.9 +
0.2011\,T$ (\unit{kJ}, $T$ em \unit{K}), válida até o ponto de fusão do zinco.
\end{solution}

\begin{solution}{exo:b2:ellingham:3}
A \qty{1500}{K}, a reta do \ce{Cr2O3} fica perto de \qty{-495}{kJ}. Abaixo dela,
portanto capazes de reduzir o óxido de cromo: silício, titânio, alumínio,
magnésio, cálcio. Acima dela, incapazes: cobre, ferro, zinco. O carbono, a
\qty{-488}{kJ}, fica logo acima: ele só reduz o óxido de cromo a uma
temperatura um pouco mais alta, e então dá um cromo carregado de carboneto
--- uma das razões pelas quais o cromo para ligas isentas de carbono é
produzido por aluminotermia.
\end{solution}

\begin{solution}{exo:b2:ellingham:4}
A inclinação é $-\Delta_r S^\circ$, governada pela variação da quantidade de
gás. \ce{C + O2 -> CO2}: um mol de gás dá um, $\Delta_r S^\circ \approx 0$,
horizontal. \ce{2C + O2 -> 2CO}: um dá dois, $\Delta_r S^\circ > 0$, a reta
desce. \ce{2CO + O2 -> 2CO2}: três dão dois, $\Delta_r S^\circ < 0$, a reta
sobe, como as dos metais.
\end{solution}

\begin{solution}{exo:b2:ellingham:5}
$\Delta_r G^\circ(\qty{600}{K}) = -181.6 + 600 \times 0.2165 = \qty{-51.7}{kJ}$,
$p_{\ce{O2}} = p^\circ\exp(-51\,710/(8.314 \times 600)) = \qty{3.2e-5}{bar}$. O ar
contém \qty{0.21}{bar}, muito mais: o mercúrio é oxidado a \qty{600}{K}
(lentamente --- foi assim que o óxido foi preparado pela primeira vez), e o
óxido só se decompõe acima de cerca de \qty{730}{K}.
\end{solution}

\begin{solution}{exo:b2:ellingham:6}
$\Delta_r G^\circ = -1582.3 - (-743.5) = \qty{-838.8}{kJ}$. Os reagentes são
dois sólidos em contato apenas na superfície de seus grãos, e a reação tem
uma energia de ativação alta: ela precisa ser iniciada localmente (um pavio
de magnésio); depois, o calor que libera a mantém em andamento.
\end{solution}

\begin{solution}{exo:b2:ellingham:7}
A reta de \ce{2C + O2 -> 2CO} cruza a de \ce{2Fe + O2 -> 2FeO} perto de
\qty{1040}{K}; acima, \ce{FeO + C -> Fe + CO} tem uma \omterm{def:b2:reaction-free-energy:reaction-gibbs}{energia de Gibbs padrão
de reação} negativa.
\end{solution}

\begin{solution}{exo:b2:ellingham:8}
$\Delta_r G^\circ = 2(-209.1) - (-396.0) = \qty{-22.2}{kJ}$, $K^\circ =
\exp(22\,200/(8.314 \times 1100)) = 11.3$. Então $x^2/(1 - x) = 11.3$, $x =
0.92$. Ao resfriar até \qty{700}{K}, o equilíbrio recua para o lado do \ce{CO2}
($x = 0.016$): o monóxido de carbono se desproporciona, \ce{2CO -> C + CO2}, e
deposita-se fuligem --- lentamente, a menos que uma superfície como a do ferro
catalise a reação.
\end{solution}

\begin{solution}{exo:b2:ellingham:9}
A reta da água (três mols de gás dão dois) sobe menos abruptamente que a
reta \ce{CO}/\ce{CO2}, também três para dois, mas com uma perda de entropia
maior. As duas se cruzam perto de \qty{1100}{K}: abaixo, a reta
\ce{CO}/\ce{CO2} é mais baixa e o monóxido de carbono é o redutor mais forte;
acima, a reta da água é mais baixa e o hidrogênio vence. É o mesmo enunciado
que o \omterm{def:b2:equilibrium-shifts:shift}{deslocamento do equilíbrio} do gás de água \ce{CO + H2O <=> CO2 + H2}
para a esquerda a alta temperatura.
\end{solution}

\begin{solution}{exo:b2:ellingham:10}
Por mol de \ce{O2}, isto é, para \ce{2MgO + Si -> SiO2 + 2Mg(g)}: $\Delta_r
G^\circ = -643.7 - (-845.5) = \qty{+201.8}{kJ}$. Com silício e sílica de
atividade 1, $\Delta_r G = \Delta_r G^\circ + RT\ln(p_{\ce{Mg}}/p^\circ)^2$, negativa
quando $p_{\ce{Mg}} < p^\circ\exp(-201\,800/(2 \times 8.314 \times 1500)) =
\qty{3e-4}{bar}$. Um forno a vácuo bombeia o vapor de magnésio à medida que
ele se forma e o condensa em uma superfície fria, de modo que sua pressão
nunca atinge esse valor; a cal adicionada à carga fixa a sílica e diminui sua
atividade, o que ajuda ainda mais.
\end{solution}

\begin{solution}{exo:b2:ellingham:11}
Quatro espécies, uma reação, três fases (dois sólidos e o gás): $v = 4 - 1 +
2 - 3 = 2$. A temperatura e pressão dadas, a razão $p_{\ce{CO}}/p_{\ce{CO2}}$
fica fixada: o gás não consegue ceder todo o seu monóxido de carbono ao óxido
de ferro, e o gás de topo ainda contém muito dele --- que é queimado para
preaquecer o ar.
\end{solution}

\begin{solution}{exo:b2:ellingham:12}
$\log K = 2(0.34 + 0.41)/0.059 = 25.4$, $K = \num{3e25}$. With
$[\ce{Fe^2+}] = \qty{1.0}{mol/L}$, $[\ce{Cu^2+}] = 1/K \approx
\qty{4e-26}{mol/L}$: o cobre é removido completamente.
\end{solution}

\begin{solution}{pb:b2:ellingham:1}
\textbf{1.} $\Delta_r H^\circ = \qty{-700.92}{kJ}$, $\Delta_r S^\circ =
\qty{-201.1}{J/K}$.
\textbf{2.} $\Delta_r G^\circ = -700.9 + 0.2011\,T$ (\unit{kJ}), abaixo de
\qty{692.7}{K}.
\textbf{3.} $\Delta_{\text{fus}}S = 7320/692.7 = \qty{10.6}{J/(K.mol)}$;
$\Delta_{\text{vap}}S = 115\,300/1180.2 = \qty{97.7}{J/(K.mol)}$.
\textbf{4.} Zinco líquido: $\Delta_r H^\circ = -700.92 - 2(7.32) = \qty{-715.6}{kJ}$,
$\Delta_r S^\circ = -201.1 - 2(10.6) = \qty{-222.2}{J/K}$: $\Delta_r G^\circ =
-715.6 + 0.2222\,T$.
\textbf{5.} Zinco gasoso: $\Delta_r H^\circ = -715.6 - 2(115.3) = \qty{-946.2}{kJ}$,
$\Delta_r S^\circ = -222.2 - 2(97.7) = \qty{-417.7}{J/K}$: $\Delta_r G^\circ =
-946.2 + 0.4177\,T$.
\textbf{6.} Inclinações $0.201$, $0.222$ e $0.418$ \unit{kJ/K}. Acima do ponto
de ebulição, a reação consome três mols de gás em vez de um: a perda de
entropia, e portanto a inclinação, dobra.
\textbf{7.} A \qty{692.7}{K}: $-700.9 + 139.3 = -715.6 + 153.9 = \qty{-561.6}{kJ}$;
a \qty{1180.2}{K}: $-715.6 + 262.3 = -946.2 + 492.9 = \qty{-453.3}{kJ}$. As
retas se unem, como devem.
\textbf{8.} Inclinação $(-453.0 + 418.2)/200 = \qty{-0.174}{kJ/K}$, $a = -418.2 +
0.174 \times 1100 = \qty{-226.8}{kJ}$: $\Delta_r G^\circ = -226.8 - 0.174\,T$.
\textbf{9.} $\Delta_r S^\circ = \qty{+174}{J/K}$: um mol de gás dá dois.
\textbf{10.} \ce{2ZnO + 2C -> 2Zn(g) + 2CO}, $\Delta_r G^\circ = (-226.8 - 0.174\,T)
- (-946.2 + 0.4177\,T) = 719.4 - 0.592\,T$ (\unit{kJ}).
\textbf{11.} Nula em $T = 719.4/0.592 = \qty{1215}{K}$.
\textbf{12.} Acima de seu ponto de ebulição, o zinco se forma como vapor: a
reação produz gás a partir de sólidos, sua entropia é grande e positiva, e o
vapor deixa a carga, o que separa o metal do minério e do coque por
destilação.
\textbf{13.} Quatro espécies, uma reação, uma condição particular ($p_{\ce{Zn}} =
p_{\ce{CO}}$), três fases: $v = 4 - 1 - 1 + 2 - 3 = 1$. Fixar a pressão
fixa a temperatura de equilíbrio.
\textbf{14.} Um mol de vapor de zinco por mol de monóxido de carbono:
$p_{\ce{Zn}} = p_{\ce{CO}} = \qty{0.5}{bar}$.
\textbf{15.} Por mol de \ce{ZnO}, $\Delta_r G^\circ = \frac12(719.4 - 0.592 \times
1300) = \qty{-25.1}{kJ/mol}$, $K^\circ = \exp(25\,100/(8.314 \times 1300)) = 10$;
$Q = 0.25 < K^\circ$: a redução ocorre.
\textbf{16.} Logo acima do cruzamento, a redução já é favorecida (com
\qty{0.5}{bar} de cada gás, a partir de cerca de \qty{1170}{K}); aquecer mais
custa combustível, desgasta as retortas de argila e não produz mais zinco,
já que a reação vai até o fim de qualquer forma.
\textbf{17.} \ce{Zn(g) + CO2 -> ZnO + CO}.
\textbf{18.} A reta do óxido de zinco fica abaixo da reta \ce{CO}/\ce{CO2}
($-445$ contra \qty{-357}{kJ} a \qty{1200}{K}): o vapor de zinco reduz o dióxido
de carbono e volta a ser óxido. Na retorta, o carbono quente destrói todo
\ce{CO2} (Boudouard); no condensador, mais frio, nada o destrói, e o \ce{CO2} ---
vindo de ar que entra por vazamentos, ou de \ce{2CO -> C + CO2} durante o
resfriamento --- estraga o zinco, que vira um pó cinzento oxidado.
\textbf{19.} Um resfriamento lento deixa o vapor por muito tempo na faixa em
que ele se reoxida e dá uma poeira fina em vez de metal líquido; uma
condensação rápida ao estado líquido limita ambos os efeitos.
\textbf{20.} $10^6/65.38 = \qty{1.53e4}{mol}$ de zinco, outros tantos de carbono:
$1.53 \times 10^4 \times 12.011 = \qty{184}{kg}$ (muito mais na prática, para
aquecer as retortas).
\textbf{21.} $\Delta_r H^\circ = \frac12(-226.8 + 946.2) = \qty{+360}{kJ}$ por
mol de zinco; por tonelada, $1.53 \times 10^4 \times 360 = \qty{5.5e6}{kJ}$, isto é,
\qty{5.5}{GJ} --- fornecidos queimando mais carvão em volta das retortas.
\textbf{22.} O carbono reduz o óxido de zinco em condições padrão acima de
$\boldsymbol{T \approx \qty{1.22e3}{K}}$.
\end{solution}
